\documentclass[10pt]{article}

\newcommand{\students}{Flavia Macovei, Arian Urdu}
\newcommand{\studentstable}{Flavia Macovei & Arian Urdu}
\newcommand{\studentsnumbertable}{7610738 & 1234567}
\newcommand{\coursetype}{Übung}

\newcommand{\fontstyle}{sans} % serif or sans
\usepackage{./01-src/head}
\usepackage{color}
\usepackage{calc}

\setcounter{exercisenumber}{1}

\begin{document}
    \maketitlepage

    \begin{abstract}
        %TODO
    \end{abstract}

    \section{Introduction}
    % ease in
	This report explores strategies for enhancing the efficiency of natural
	language processing under constraints of limited training data and
	computational resources. Traditionally, “bigger” language models, i.e. models
	with more parameters produce better predictions than their smaller counterparts.
	This approach of scaling up models usually involves large datasets and powerful
	hardware, both of which can be expensive. % <- umformulieren + zitat
	In this report, we consider alternative methods for improving the efficiency
	of a model and aim to maximise prediction quality with strict constraints on
	dataset size and hardware power.

	More specifically, the goal of this project
	is to develop a small language model that is trained on a curated,
	small-scale dataset and whose training process can be conducted on a
	standard laptop. The limited size of the dataset entails a need for
	specialisation of the prediction domain, so we aim to achieve high
	prediction accuracy in a specific target area. The model is to be designed in a way that it takes an input prompt in text form and predict its continuation.

    Eldan and Li pose a similar question to this report's objective in their paper which first introduced the TinyStories dataset: “Language models ... often struggle to produce coherent and fluent text when they are small. ... This raises the question of whether the emergence of the ability to produce coherent English text only occurs at larger scales.” % <- zitat
    The TinyStories dataset consists of roughly 2 million short stories generated with GPT-3.5 and GPT-4 in simple English and lends itself to be used in the scope of this project. Eldan and Li argue that natural language prediction has to meet two criteria to be considered coherent: it should fit both content and the syntax of its context.

    The first suggests a form of logical reasoning ability of the model. Take for an example the sentence “Jessica was cold so she put on ...”. % <- schöner formattieren
    The expected completion generated by a language model should relate to a warm item of clothing. The second implies the ability to use proper grammar: not any phrase that relates to clothing can be used as a continuation of the above prompt, but only those that adhere to grammatical and syntactical rules of the English language.

    Building on these criteria, the objectives of this report are as follows:

    \subsection{Reproducing stories}
    The model should be able to generate stories similar to those of its training set given a short prompt or no prompt.
    They should resemble the training stories in length, language and plot complexity.
    A model which achieves this objective demonstrates the ability to string words together coherently, fulfilling the grammatical and syntactical requirement.
    Additionally, it is desired that the model display a certain degree of creativity and not directly replicate its training data.

    \subsection{Logical inference}
    Secondly, a successful model is expected to be able to generate a prediction that fits the logic and content of its context. This comes into play when the input prompt is longer and provides more context, thereby restricting the realm of feasible predictions.
    This objective feeds into the requirement for logical conformity mentioned above.

    \subsection{Competitive with leading models}
    Finally, the model outlined in this project is desired to demonstrate competitive performance when compared to leading benchmarks.
    Specifically, it should perform as well as or better than those models on selected evaluation metrics. This comparison is an indication of the model's overall capability and of the validity of the methods presented in this report.

    \section{The Dataset}

    - repeat reasons for choosing it (very short)\\
    - origin\\
    - facts\\

    \section{The Model}
    - parameters\\
    - performance\\
    \quad > evaluation metrics + results\\
    \quad > argue for chosen metrics\\
    \quad > (running time)\\

    \section{Results}
    1. preprocessing\\
    2. tokenisers\\
    3. curriculum learning + layer freezing\\

    \section{Future Work}

    \section{Conclusion}

\end{document}